% Part: computability
% Chapter: computability-theory
% Section: ce-sets

\documentclass[../../../include/open-logic-section]{subfiles}

\begin{document}

\olfileid[tr]{cmp}{thy}{ces}
\olsection{Hesaplanabilir Biçimde Numaralandırılabilir Kümeler}

\begin{defn}
Bir küme boşsa ya da hesaplanabilir bir fonksiyonun görüntü kümesiyse
\emph{hesaplanabilir biçimde numaralandırılabilir}dir.
\end{defn}

\begin{history}
Hesaplanabilir biçimde numaralandırılabilir kümelere \emph{özyinelemeli
biçimde numaralandırılabilir} de denir. Bu özgün terminolojidir; günümüzde
iki terim ile ``c.e.'' ve ``r.e.'' kısaltmalarının tümü yaygın biçimde
kullanılır.
\end{history}

\begin{explain}
Tanımın ne anlama geldiğini ve terminolojinin neden uygun olduğunu düşünün.
Düşünce şudur: $S$, hesaplanabilir $f$ fonksiyonunun görüntü kümesiyse
\[
  S=\{f(0),f(1),f(2),\dots\}
\]
olur ve $f$, $S$ elemanlarını ``numaralandırıyor'' diye görülebilir. Tanıma
göre $f$'nin artan bir fonksiyon olmasının gerekmediğine, yani numaralandırma
sırasının artan olması gerekmediğine dikkat edin. Hatta $f$'nin bire bir
olması bile gerekmez; $S$'nin $f(0),f(1),f(2),\dots$ numaralandırmasında
tekrarlara izin verilir. Örneğin sabit $f(x)=0$ fonksiyonu $\{0\}$ kümesini
numaralandırır.
\end{explain}

Her hesaplanabilir küme hesaplanabilir biçimde numaralandırılabilir. Bunu
görmek için $S$ hesaplanabilir olsun. $S$ boşsa tanım gereği hesaplanabilir
biçimde numaralandırılabilir. Değilse $a$, $S$'nin herhangi bir elemanı olsun
ve $f$'yi
\[
  f(x)=\begin{cases}
    x & \text{$\Char{S}(x)=1$ ise}\\
    a & \text{aksi durumda}
  \end{cases}
\]
ile tanımlayalım. O zaman $f$ hesaplanabilir ve $S$, $f$'nin görüntü
kümesidir.

\end{document}

